% !TeX root = ../papers/nc-p3.tex
\section{混合交換境界の全特殊値の正則性}
\label{sec:mixed-exchange-regularity}
\begin{theorem}\label{thm:mixed-exchange-regularity}
定理\ref{thm:klein-mixed-exchange-resolution}の $z=1$, $\lambda\ne0,\pm1$ の
表示で、$x=ad-fg$, $\det A$, $\det B$ を非零とする。
標準Koszul左右AS正則性の必要十分条件は
\[
                 k=(tb-\lambda)(tb-1/\lambda)\ne0
\]
である。$t=0$ の全点を含む。$k=0$ の全点は $4n$ 文字の商空間に
少なくとも $4^n$ の次元を持つ始点を有する。
\end{theorem}
\begin{proof}
規格化前の二つの式から
$(ta-1)(td-1)=(tf-1)(tg-1)=k$ を得る。
$t\ne0$ では $a,b,d,f,g$ の共通倍率で $t=1$ にできる。
$k\ne0$ のとき $(a,d,f,g)=(1+u,1+k/u,1+v,1+k/v)$ であり、
$x=(u-v)(1-k/(uv))$。$t=0$ では $a+d=f+g=(\lambda+1/\lambda)b$ である。

まず $k\ne0$ とする。偶数辺の薄い二関係は全係数1である。
最初の偶数辺で厚色の係数は前節の $D_0,D_1$ のグラフをなす。
$Y_\lambda=\lambda P X_\lambda$ と書く（$P$ は二文字の置換）。
$D_0=Y_\lambda A^{-1}$ の全成分は非零である。
実際、零なら $a,d$ のいずれかが $\lambda b,b/\lambda$ のいずれかに等しく、
$t\det A=a+d-(\lambda+1/\lambda)b$ と $\det A\ne0$ により $k=0$ となる。
階数1関係があるとする。薄い二係数の符号付き和の一方が零でなければ
交差ブロックは可逆であり不可能。薄い二係数がともに零なら厚色のグラフ性に反する。
従って二つの交差ブロックはともに $(1,\sigma)^{\mathsf t}(1,\sigma)$ の
非零倍である（$\sigma=\pm1$）。二つの対角ブロックも同じ行列の非零倍となる。
これは $D_0,D_1$ が $(1,1)^{\mathsf t}$ を同じ倍数で保つことを要する。
$A,B^{\mathsf t}$ がこの直線を保つので $a=d,f=g$、さらにその固有値の一致から
$(1+\lambda)/(a+b)=(1+\lambda)/(b+f)$、従って $x=0$ で矛盾する。

同じ偶数辺の左評価が低階数になるとする。前節と同じ射影で両半分の支持数は等しい。
二支持では、二つの薄い関係が定めるグラフは単項行列である。色0の評価をこの
グラフに入れると $D_0$ の一成分が零になり矛盾する。
全支持では $p=(1,s,1,m)$ とする。前節の薄い二行列が可逆なら、$z=1$ とした
第二型の式から $B=PA^{\mathsf t}$ が従う（$T$ が $P$ と可換であることでも確認できる）。
従って $C$ は特異となる。薄い二行列が特異なら $s=m=\pm1$ である。
厚い二色の式から $D_1=PD_0P$、ゆえに $B=A^{\mathsf t}P$ となり、やはり $x=0$。
反対の偶数辺では $A,B,\lambda$ を $A,B^{\mathsf t},1/\lambda$ に取り替えると
同じグラフになる。右評価には辺の関係行列を転置すればよく、色1のグラフは
$PD_1P$ となる。これは $B$ と $PBP=B^{\mathsf t}$ の交換である。
以上の操作は $x$, $k$ と非消滅条件を保つ。

奇数辺では、厚い二色の係数平面がそれぞれ
\[
 (u_0,u_1,\alpha u_0+\beta u_1,\alpha u_0+\beta u_1),\qquad
 (\gamma v_0+\delta v_1,\gamma v_0+\delta v_1,v_0,v_1)
\]
となり、四つの係数は全て非零である。色0の基底は $\binom{C}{J}$ または
$\binom{C^{\mathsf t}}{J}$、色1では上が $J$、下が可逆な角である。$J$ は全成分1の行列とする。
\path{computations/verify_mixed_exchange_regularity.py} は、上記の $t=0,1$ の
両座標で残る角と各係数を計算し、その分子・分母の既約因子が
$\lambda(\lambda^2-1)kx\det A\det B$ と座標の単元の因子であることを検算する。
奇数辺の薄い色の二つの交差行列式は比例しない（積比は $\lambda^{\pm2}\ne1$）。
よって純関係では薄い係数がともに零となり、一つの対角ブロックが零になる。
残るブロックには一色の一次元しか残らず、その二つの成分はともに非零で階数2である。

奇数辺の評価点で低位半分の両座標が非零なら、色0の像は低位への射影が同型の
二平面で、高位での像は高位座標対が張る直線以下である。高位座標が非零なら
色1の高位像はこの直線から出る（支持二なら全射、支持一なら別の座標）。
高位の両座標が非零の場合も役割を逆にして同じである。片半分だけの支持なら
厚い二色でその半分を張り、薄い色で別の方向を得る。
一座標ずつの支持なら薄い二色だけで両半分間の単項同型のグラフが得られ、
色0がそこに入るには $\alpha\beta=0$ が必要である。
右評価では各色内の添字が入れ替わるだけで結論は同じである。
一文字全階数は薄い非零成分と文字次数分解から従う。
定理\ref{thm:involution-quartet-component}と混合表示の特殊化定理を適用する。

$k=0$ なら $t=1$ にでき、$b=\lambda,1/\lambda$ である。
$a,d$ の第 $i$ 成分、$f,g$ の第 $j$ 成分を1に選ぶ（$i,j=0,1$）。
残りを $u,v$ とし、$b=\lambda,1/\lambda$ の順に $e=0,1$ とする。
$i=j$ なら位置 $1+e$、$i\ne j$ なら位置 $3e$ に
$p=e_i+e_{i+2}$ を置く。両隣の核平面と中央三文字商を前節と同じ方法で取る。
検算では奇数辺の色1だけは最後の二列を基底に使う。核平面の薄い二行の低位小行列式は
単項式である。八表示の全てで二平面が六関係を消し、三独立変数
$u,v,\lambda$ 上で挿入関係の階数は12である。特殊化で上がらないため
三文字商の次元は4以上となり、命題\ref{prop:cut-arrow-growth}で結論を得る。
\end{proof}

\begin{corollary}\label{cor:regular-full-thin-all-pairings}
二隣接切断の色次元が $(2,2,1,1)$ の並べ替えで四薄色の全成分が非零な四元群窓では、
候補の正則性に全パラメータで有限な必要十分判定がある。
\end{corollary}
\begin{proof}
異方向は定理\ref{thm:klein-transverse-dense}を使う。
平行方向の被覆は定理\ref{thm:full-thin-parallel-cover}である。両交換境界の
正負表示は定理\ref{thm:full-thin-double-exchange}、一交換境界は直前の定理で決まる。
残った可逆交差部の主枝には定理\ref{thm:elliptic-evaluation-obstructions}を使う。
二つの交換枝には定理\ref{thm:bisymmetric-two-length-test}を正表示に使い、
負表示は指定コサイクルを外して同じ検査をする。文字次数を保つ捻りはKoszul性、
左右の完全対合、各始点のHilbert関数を保つ。各閉包への所属も、明示された
有限型表示の稠密開部のグラフ閉包を射影的に消去して判定できる。
\end{proof}
